\chapter{预测性补充分析}

第四章已经通过回归模型检验了 AI 文旅短视频内容特征对传播效果的影响。本章进一步引入机器学习方法，从预测视角对前文结论进行补充验证。之所以将该部分单独成章，而不是放入第三章研究设计，是因为第三章主要回答“如何提出解释模型和假设”，第四章回答“哪些变量具有统计解释力”，本章则回答“同一套变量能否在预测任务中仍然提供信息”。因此，本章属于结果后的交叉验证环节，并不替代回归模型的假设检验功能，而是用于考察现有变量体系是否具有预测能力，非线性模型是否能够提供额外信息，以及核心变量在预测任务中是否仍然保持较高重要性。

\section{分析定位}

本章引入机器学习方法，目的在于为前文回归结果提供另一种验证视角。回归分析侧重解释变量影响方向和显著性，能够说明某一因素在控制其他变量后是否仍然发挥作用。短视频传播效果则具有更强的复杂性，互动量并不一定由单一变量线性决定。AI 技术质感、视频风格、视频时长和发布时间等因素可能共同影响用户的观看体验和互动意愿。例如，技术质感较高的视频如果同时具有鲜明风格和较紧凑的视频节奏，可能比只具备其中某一项优势的视频更容易获得互动。

机器学习方法适合从整体预测角度检验变量体系的有效性。如果同一套变量不仅能够在回归模型中解释传播效果，也能在预测模型中区分互动表现，则说明这些变量具有较强稳定性。此时，机器学习并不是为了追求更复杂的模型形式，而是用于观察内容变量在预测任务中是否仍然具有实际信息量。

本章还将连续互动量预测扩展为高互动视频识别。地方文旅账号在实际运营中并不只关心某个变量是否显著，也关心哪些内容特征更可能帮助视频进入高互动区间。通过将总互动量前 25\% 的样本定义为高互动视频，可以进一步观察现有变量体系是否能够识别潜在高互动内容。

因此，本章的核心目的，是在不改变样本口径和变量体系的前提下，从预测视角进一步检验 AI 技术质感、视频风格、视频时长等核心变量对传播效果的解释价值。

\section{目标与特征}

本章继续使用前文完成深度编码的 722 条 AI 文旅短视频样本。样本来源于抖音平台省级文旅账号发布的 AI 相关短视频，并经过筛选、复核和人工内容编码后形成最终分析数据。为了保证结果具有可比性，本章不重新引入新的样本，而是沿用第四章回归分析使用的同一套变量体系。

传播效果的目标变量为总互动量。本文将点赞量、评论量、收藏量和转发量相加，得到每条视频的总互动量。短视频互动数据通常呈现明显长尾分布，少数视频互动量极高，直接使用原始互动量容易受到极端值影响。为降低长尾分布对模型训练的干扰，本章将目标变量设定为 $y_i=\log(\text{总互动量}_i+1)$，其中 $y_i$ 表示第 $i$ 条视频的机器学习预测目标。

自变量主要来自前文内容编码结果，包括 AI 介入形式、AI 介入程度、AI 技术质感、视频主体、视频主题、内容导向、视频风格、视频时长和发布距采集日天数。为了观察地区差异是否会补充解释传播效果，本章设置两类特征集。一类只包含主体变量，主要反映视频内容和 AI 呈现特征；另一类在主体变量基础上加入省份变量，用于观察地区文旅资源、账号基础和地方传播环境可能带来的影响。

在数据预处理方面，分类变量转化为哑变量，连续变量采用对数化后的口径，并在建模过程中进行标准化处理。模型评价采用 5 折交叉验证，以减少单次训练集和测试集划分带来的偶然影响。

\section{模型选择}

本章连续型传播效果预测使用 Ridge 线性基准模型、随机森林模型和梯度提升模型。三类模型的复杂度逐步提高，分别承担线性基准比较、非线性关系识别和更强预测模型检验的功能。本文样本量为 722 条，变量主要来自人工内容编码，因此模型选择并不追求复杂算法堆叠，而是优先考虑模型含义清晰、结果便于解释、能够与前文回归分析形成衔接的基础机器学习方法。这样的安排能够回答三个层次的问题：线性模型是否已经能够利用这些变量，非线性模型是否提供额外增益，以及复杂模型的提升是否足以改变前文判断。

Ridge 线性基准模型是在普通线性回归基础上加入 L2 惩罚项的模型。Hoerl 和 Kennard 提出的岭回归方法通过向 $X^{\mathrm{T}}X$ 的对角线上加入正值，缓解解释变量非正交时最小二乘估计容易出现的系数不稳定问题\cite{hoerl1970ridge}。在本文样本中，分类变量和哑变量较多，部分内容特征之间也可能存在相关性，因此 Ridge 模型能够在保留线性解释框架的同时降低过拟合风险。其目标函数可表示为
\begin{compactformula}
\[
\hat{\beta}=\arg\min_{\beta}\left[\sum_{i=1}^{n}(y_i-\mathbf{x}_i^{\mathrm{T}}\beta)^2+\lambda\sum_{j=1}^{p}\beta_j^2\right].
\]
\end{compactformula}
式中，$\mathbf{x}_i$ 表示第 $i$ 条视频的解释变量向量，$\lambda$ 为惩罚系数。该模型在本章中的作用是提供线性预测基准。若 Ridge 模型已经具有一定预测能力，说明本文构建的内容变量本身包含有效信息；若后续非线性模型进一步提高预测表现，则说明变量之间可能存在普通线性模型难以充分捕捉的组合关系。由于 Ridge 模型与第四章回归分析在形式上较为接近，它也有助于保持机器学习结果与前文解释性分析之间的可比性。

随机森林模型是一种基于多棵决策树的集成学习方法。Breiman 提出的随机森林将多棵树预测器组合起来，使每棵树在样本抽取和变量选择上都具有一定随机性，并通过集成结果降低单棵树对样本划分的敏感性\cite{breiman2001random}。单棵决策树会根据变量取值不断划分样本，形成一组判断规则，但其结果容易受到局部划分的影响。随机森林通过重复抽取样本、随机选择部分变量并训练多棵树，再对多棵树的预测结果取平均，从而提高预测稳定性。其预测结果可表示为
\begin{compactformula}
\[
\hat{f}_{RF}(\mathbf{x}_i)=\frac{1}{B}\sum_{b=1}^{B}T_b(\mathbf{x}_i).
\]
\end{compactformula}
其中，$T_b(\mathbf{x}_i)$ 为第 $b$ 棵回归树的预测结果，$B$ 为树的数量。随机森林不要求变量与目标变量之间完全符合线性关系，能够捕捉阈值效应和变量组合关系。对于本文而言，该模型适合观察 AI 技术质感、视频风格、视频时长等内容变量是否在组合后产生更强的传播区分度。例如，高质感画面与强视觉风格同时出现时，模型可以通过树结构识别这种联合特征，而不必预先设定具体交互项。随机森林还能够输出变量重要性，因此本章后续变量重要性分析基于该模型展开。

梯度提升模型同样属于树模型，但其建模思路不同于随机森林。Friedman 将 boosting 方法解释为函数空间中的逐步加性优化过程，模型每一轮都沿着损失函数下降方向加入新的弱学习器，从而不断修正已有预测结果\cite{friedman2001greedy}。随机森林侧重并行生成多棵树并进行平均，梯度提升则强调逐步修正误差。模型先建立一个初始预测结果，再不断加入新的弱学习器，使新模型尽量弥补前一轮预测中的不足。其预测函数可写为
\begin{compactformula}
\[
\hat{f}_{M}(\mathbf{x}_i)=\hat{f}_{0}(\mathbf{x}_i)+\sum_{m=1}^{M}\eta h_m(\mathbf{x}_i).
\]
\end{compactformula}
其中，$h_m(\mathbf{x}_i)$ 为第 $m$ 轮加入的弱学习器，$\eta$ 为学习率，$M$ 为迭代轮数。每一轮模型都在前一轮结果基础上继续改善预测表现，因此更适合检验当前变量体系在较强非线性模型下是否仍有提升空间。若梯度提升模型明显优于 Ridge 和随机森林，说明传播效果可能存在更复杂的非线性结构；若提升幅度有限，则说明当前人工编码变量的预测能力存在边界，继续增加模型复杂度未必能够显著改善解释效果。

\section{检验步骤}

本章检验路径由连续型传播效果预测、变量重要性分析和高互动视频识别三部分组成。三部分之间并非彼此割裂，而是依次回答“能否预测”“依赖哪些变量”以及“能否识别高互动内容”。

连续型传播效果预测以 $\log(\text{总互动量}+1)$ 为目标变量，分别使用 Ridge 线性基准模型、随机森林模型和梯度提升模型进行建模。实验在两类特征集上展开：主体变量用于观察视频内容和 AI 呈现特征本身是否具有预测价值，加入省份变量后的特征集则用于检验地区因素是否提供额外信息。模型表现主要通过 $R^2$、RMSE 和 MAE 进行比较。$R^2$ 越高，说明模型对互动表现的解释和预测能力越强；RMSE 与 MAE 越低，说明预测误差越小。

变量重要性分析基于随机森林回归模型展开。变量重要性不同于回归系数，不能直接说明变量影响方向，也不能替代显著性检验。它反映的是模型在预测互动表现时对某一变量的依赖程度。如果打乱某一变量后模型预测表现明显下降，说明该变量在预测任务中具有较高贡献。

高互动视频识别将连续的互动量结果转化为二分类问题。本文把总互动量位于样本前 25\% 的视频定义为高互动视频，其余视频定义为普通互动视频，并使用随机森林分类器进行识别。该部分不是新增一个与前三个预测模型并列的主模型，而是从运营实践角度补充回答现有变量体系能否识别更可能进入头部互动区间的视频。

\section{预测结果}

三个模型对连续型传播效果的预测结果见表 \ref{tab:mlresult}。

\begin{table}[!htbp]
  \centering
  \caption{机器学习模型预测效果比较}
  \label{tab:mlresult}
  \small
  \begin{tabular}{llccc}
    \toprule
    特征集 & 模型 & $R^2$ 均值 & RMSE 均值 & MAE 均值 \\
    \midrule
    主体变量 & Ridge 线性基准 & 0.457 & 1.442 & 1.127 \\
    主体变量 & 随机森林 & 0.477 & 1.415 & 1.078 \\
    主体变量 & 梯度提升 & 0.481 & 1.412 & 1.097 \\
    加入省份 & Ridge 线性基准 & 0.559 & 1.300 & 0.988 \\
    加入省份 & 随机森林 & 0.505 & 1.378 & 1.047 \\
    加入省份 & 梯度提升 & 0.561 & 1.298 & 1.003 \\
    \bottomrule
  \end{tabular}
\end{table}

在仅使用主体变量时，Ridge 线性基准模型的 $R^2$ 为 0.457，说明视频内容与 AI 呈现变量已经能够解释一部分传播效果差异。随机森林模型和梯度提升模型的 $R^2$ 分别为 0.477 和 0.481，略高于 Ridge 线性基准模型。这一结果表明，非线性模型能够带来一定提升，传播效果并非完全由单一变量线性决定，而是可能存在一定组合关系。

加入省份变量后，Ridge 线性基准模型和梯度提升模型的 $R^2$ 分别提高至 0.559 和 0.561，RMSE 与 MAE 也有所下降。省份变量带来的增益说明，地区文旅资源基础、账号粉丝规模、地方传播环境和长期运营能力等因素可能会影响互动表现。视频内容变量具有重要价值，但传播效果并不只由视频内容决定，地区和账号基础同样构成了传播结果的重要背景。

复杂模型并没有取得压倒性优势。梯度提升模型整体表现最好，但相较随机森林和 Ridge 模型的提升幅度有限。由此可见，本章机器学习结果更适合作为回归结论的补充验证，而不宜被解释为能够精确预测单条视频流量的模型。

变量重要性结果见图 \ref{fig:rf}。

\begin{figure}[!htbp]
  \centering
  \includegraphics[width=0.92\textwidth]{../analysis_outputs/ml_experiment/02_随机森林变量重要性.png}
  \caption{随机森林变量重要性}
  \label{fig:rf}
\end{figure}

从图 \ref{fig:rf} 可以看出，AI 技术质感、视频时长、视频风格和发布距采集日天数是较重要的变量。其中，AI 技术质感的重要性最高，说明模型在预测互动表现时最依赖视频最终呈现出来的画面质量和视觉完成度。用户并不是因为视频简单贴上了 AI 标签就产生互动，而是更容易对更精良、更有完成度的 AI 视觉表达作出反应。

视频时长的重要性说明，短视频平台中的观看成本对传播效果具有约束作用。过长或节奏不紧凑的视频可能降低用户完整观看和进一步互动的可能性。视频风格的重要性则表明，抽象、超现实或具有强视觉识别度的表达方式，有助于提升内容在信息流中的辨识度。发布距采集日天数的重要性提醒我们，公开互动数据具有时间累积特征，较早发布的视频通常拥有更长的互动积累窗口。

高互动识别实验结果显示，随机森林分类器的 AUC 为 0.851，F1 值为 0.646，说明当前变量体系对高互动视频具有一定区分能力。AUC 主要衡量模型区分高互动视频和普通互动视频的能力，F1 值综合考虑准确率和召回率。该结果表明，AI 技术质感、视频风格和视频时长等变量不仅能够解释整体互动量差异，也有助于识别更可能进入头部互动区间的视频。

机器学习结果与第四章回归分析在核心判断上具有连续性。回归分析提供变量方向和显著性，机器学习分析提供整体预测能力、变量重要性排序和高互动识别表现。两类方法都指向 AI 技术质感、视频风格和视频时长，说明这些变量在解释和预测两个视角下都具有较高稳定性。

\section{本章小结}

本章围绕 AI 文旅短视频传播效果的预测性补充验证展开，在前文 722 条深度编码样本基础上，以 $\log(\text{总互动量}+1)$ 作为预测目标，使用 Ridge 线性基准模型、随机森林模型和梯度提升模型进行连续型传播效果预测，并进一步分析变量重要性和高互动视频识别表现。

实验结果表明，本文编码的内容变量具有一定预测价值。主体变量条件下，随机森林和梯度提升模型略优于 Ridge 线性基准模型，说明变量之间可能存在一定非线性和组合效应；加入省份变量后，模型表现进一步提高，表明地区和账号层面因素也会影响传播效果。变量重要性和高互动识别结果进一步显示，AI 技术质感、视频风格和视频时长不仅在回归分析中具有解释价值，在预测任务和高互动识别中也具有较高区分度。

总体来看，机器学习分析并没有改变本文关于 AI 文旅短视频传播机制的基本判断，而是从预测视角强化了前文结论。AI 技术本身并不必然转化为更高互动，关键在于技术能否形成用户可感知的画面完成度、风格辨识度和平台节奏适配度。对于文旅账号而言，AI 内容生产的重点不应停留在技术使用层面，而应进一步落实到视觉质量控制、风格表达设计和短视频观看体验优化之中。
